\section{Costo di calcolo, alternative considerate e lavori correlati}\label{sec:cos}

Questo capitolo dà il costo di calcolo del sistema e delle rappresentazioni dello spazio con cui è stato confrontato, elenca le alternative considerate ed escluse con il motivo tecnico e con il controesempio o la misura che le esclude, e colloca il lavoro rispetto alla letteratura. Tutte le cifre di tempo e di memoria sono misure fatte su un calcolatore portatile, su profondità sintetica o su una registrazione pubblica: nessuna viene da un sensore fisico né dalla scheda di calcolo di bordo. La scelta della rappresentazione è decisa dalla sicurezza dello stato «libero» (\S\ref{sec:map}) prima che dal costo.

\subsection{Rappresentazioni dello spazio e loro complessità}\label{sec:cos-rappresentazioni}

Simboli aggiunti alla notazione comune: $W\times H$ pixel dell'immagine di profondità; $k$ passo di sottocampionamento dei pixel; $N_r\le WH/k^2$ raggi (pixel validi letti) per fotogramma; $L_x\times L_z\times L_y$ regione rappresentata (pianta e altezza); $s$ lato del voxel; $n_x,n_z,n_y$ numero di voxel per lato; $N_v=n_xn_yn_z$; $c$ byte per voxel; $\beta$ lato di un blocco, in voxel; $B$ blocchi allocati; $K$ passi di raggio per fotogramma (voxel attraversati, sommati sui raggi); $N_c$ celle della griglia polare; $n_T$ oggetti seguiti in un ciclo; $T=\Delta=1/15\unita{s}=66{,}7\unita{ms}$ periodo del ciclo.

Sono state considerate sette rappresentazioni. (i)~La \emph{griglia polare del suolo a quattro stati} (\S\ref{sec:map}): celle di $1^\circ\times0{,}1\unita{m}$ fino a $3\unita{m}$, $N_c=87\cdot30=2\,610$ con campo orizzontale di $87^\circ$; ogni pixel valido è assegnato alla sua cella con tabelle precalcolate e nessun raggio viene percorso. (ii)~La \emph{lettura diretta nell'immagine di profondità}: nessuna mappa; lo stato di un punto si legge per proiezione sugli ultimi fotogrammi (\S\ref{sec:map}). (iii)~La \emph{griglia densa di voxel}: ogni raggio è percorso voxel per voxel fino al ritorno. (iv)~I \emph{voxel sparsi con tabella hash a blocchi} di $\beta^3=8^3$ voxel, allocati alla prima scrittura, schema proposto per la fusione volumetrica in tempo reale su hardware grafico comune \cite{niessner2013}. (v)~L'\emph{octree}, di profondità 6 a $s=10\unita{cm}$ e 7 a $5\unita{cm}$, 8 byte per nodo, con o senza fusione dei figli uguali. (vi)~La \emph{mappa a fasce di quota}: griglia in pianta con tre fasce per colonna (0,15--0,60 m, 0,60--1,60 m, 1,60--2,0 m), e la variante a colonne di bit (30 livelli da 10 cm, 12 byte per colonna). (vii)~La \emph{griglia di occupazione bayesiana}, che è il contenuto di un voxel (logaritmo delle quote di occupazione) e si appoggia a una delle tre strutture precedenti. La tabella~\ref{tab:cos-complessita} ne riassume la complessità.

\begin{table}[H]\centering\small
\caption{Complessità delle rappresentazioni considerate.}\label{tab:cos-complessita}
\begin{tabular}{@{}llll@{}}\toprule
Rappresentazione & Memoria & Aggiornamento per fotogramma & Lettura di un punto\\\midrule
Griglia polare del suolo & $O(N_c)$ & $O(N_r+N_c)$ & $O(1)$\\
Lettura nell'immagine & $O(WH)$ per fotogramma tenuto & $O(WH)$ & $O(1)$: nove pixel\\
Griglia densa di voxel & $c\,N_v=c\,L_xL_yL_z/s^3$ & $O(K)$, $K=O(N_rL/s)$ & $O(1)$\\
Voxel sparsi con hash & $c\,\beta^3B$ più la tabella & $O(K)$ atteso, costante maggiore & $O(1)$ atteso\\
Octree & 8 byte per nodo & $O(K\log_2 n)$ & $O(\log_2 n)$\\
Fasce di quota & $O(n_xn_z)$ per fascia & $O(K)$ & $O(1)$\\
Occupazione bayesiana & quella del supporto & quella del supporto & $O(1)$\\\bottomrule
\end{tabular}
\end{table}

\begin{lemma}[Numero di voxel e memoria]\label{lem:cos-voxel}
Per la regione $6\times4\times3\unita{m}$ si ha $N_v=60\cdot40\cdot30=72\,000$ a $s=10\unita{cm}$ e $N_v=120\cdot80\cdot60=576\,000$ a $s=5\unita{cm}$; la memoria densa $c\,N_v$ vale $288\,000$ e $2\,304\,000$ byte con $c=4$ byte, il doppio con $c=8$ byte. I blocchi da $8^3$ voxel sono al più $160$ e $1\,200$. \prova{C}
\end{lemma}

\begin{lemma}[Tetto dei passi di raggio]\label{lem:cos-passi}
In una griglia di $n_x\times n_y\times n_z$ voxel ogni raggio attraversa al più $n_x+n_y+n_z-2$ voxel; quindi $K\le N_r\,(n_x+n_y+n_z-2)$, cioè $128\,N_r$ a $10\unita{cm}$ e $258\,N_r$ a $5\unita{cm}$ nella regione del lemma~\ref{lem:cos-voxel}, qualunque sia il numero di oggetti nella scena. \prova{D}
\end{lemma}
\begin{proof}
Passando da un voxel al successivo il raggio attraversa esattamente una superficie di separazione della griglia. Un segmento di retta attraversa ciascuna di queste superfici al più una volta, e quelle interne alla regione sono $n_x-1$, $n_y-1$ e $n_z-1$ nelle tre direzioni: i voxel toccati sono al più $1+(n_x-1)+(n_y-1)+(n_z-1)$. Un oggetto in più può solo accorciare un raggio.
\end{proof}

Nelle misure del \S\ref{sec:cos-misure} i passi per raggio sono in media 46,8 a 10 cm e 94,1 a 5 cm: dimezzando $s$ il costo raddoppia, mentre $N_v$ cresce di otto volte. Il caso peggiore per il tempo è la scena vuota, non quella affollata.

\begin{lemma}[Limite del guadagno di una struttura sparsa a tre stati]\label{lem:cos-sparsa}
Se «non allocato» deve valere «non osservato» (lo impone il controesempio della lettura a due stati, \S\ref{sec:map}), ogni voxel dichiarabile libero va allocato, e la memoria cresce con il volume visto libero, non con le superfici. Entro una portata $R$ il sensore vede al più la piramide
\begin{equation}\label{eq:cos-piramide}
V_p=\tfrac43\,\tan\!\Big(\tfrac{87^\circ}{2}\Big)\tan\!\Big(\tfrac{58^\circ}{2}\Big)R^3=\tfrac43\cdot0{,}9490\cdot0{,}5543\cdot27\unita{m^3}=18{,}9\unita{m^3}\qquad(R=3\unita{m}),
\end{equation}
cioè il 26,3\,\% dei $72\unita{m^3}$ della regione: nel caso peggiore, con tutta la piramide vista libera, il rapporto fra memoria densa e memoria sparsa è $72/18{,}9=3{,}8$, prima del costo dell'indice; è il guadagno su cui si può contare, perché con ostacoli in scena il volume visto libero è minore e il rapporto maggiore. \prova{C}
\end{lemma}

\subsection{Misure di tempo e memoria}\label{sec:cos-misure}

\textbf{Condizioni di misura.} Calcolatore portatile con processore AMD Ryzen AI 9 HX 370 (quattro nuclei fino a 5,16 GHz, otto fino a 3,29 GHz), Linux, programmi in C compilati con gcc 16.2.1 e \texttt{-O2}, un solo filo d'esecuzione fissato su un nucleo. Profondità \emph{sintetica} $640\times480$, campo $87^\circ\times58^\circ$, $h=1{,}2\unita{m}$, $\theta=15^\circ$, valida fino a $6\unita{m}$, rumore $\sigma_z=0{,}0025\,d^2$ (metri), 1,2\,\% di pixel non validi; regione $6\times4\times3\unita{m}$, $c=4$ byte; scena con parete di fondo, tre scaffali e venti oggetti; $N=300$ fotogrammi per configurazione. Il tempo comprende deproiezione, percorso dei raggi e scritture, non la generazione dell'immagine. Con $k=4$ i raggi validi sono $16\,557$, con $k=1$ sono $264\,709$.

\begin{table}[H]\centering\footnotesize
\caption{Memoria e tempo per fotogramma (media / 95° percentile, in ms) sul nucleo veloce; profondità sintetica, $N=300$ fotogrammi per casella. Per le strutture sparse la memoria va dalla scena con il solo suolo a quella ingombra. La griglia polare e la lettura nell'immagine non dipendono da $s$: le due colonne sono due esecuzioni. \prova{S}}\label{tab:cos-tempi}
\begin{tabularx}{\textwidth}{@{}Lrrcccc@{}}\toprule
 & \multicolumn{2}{c}{Memoria (byte)} & \multicolumn{2}{c}{$k=4$} & \multicolumn{2}{c}{$k=1$ (tutti i raggi)}\\\cmidrule(lr){2-3}\cmidrule(lr){4-5}\cmidrule(l){6-7}
 & 10 cm & 5 cm & 10 cm & 5 cm & 10 cm & 5 cm\\\midrule
Griglia polare, 3 m & \multicolumn{2}{c}{$10\,788$ (più 1,5 MB di tabelle)} & 0,03 / 0,04 & uguale & 0,42 / 0,55 & 0,37 / 0,40\\
Lettura nell'immagine & \multicolumn{2}{c}{$2\,457\,600$ (due fotogrammi)} & 0,51 / 0,52 & 0,51 / 0,52 & 0,74 / 0,88 & 0,51 / 0,53\\
Colonne di bit & $28\,800$ & $115\,200$ & 1,30 / 1,34 & 2,47 / 2,55 & 19,5 / 20,0 & 37,4 / 38,4\\
Tre fasce di quota & $48\,000$ & $192\,000$ & 3,29 / 3,37 & 6,23 / 6,41 & 44,5 / 45,7 & 85,5 / 97,6\\
Griglia densa & $288\,000$ & $2\,304\,000$ & 2,03 / 2,08 & 4,21 / 4,35 & 22,4 / 22,8 & 44,5 / 49,1\\
Hash a blocchi $8^3$ & $159\,744$--$241\,664$ & $763\,904$--$1\,515\,520$ & 3,02 / 3,12 & 6,71 / 7,45 & 32,6 / 47,9 & 72,4 / 76,4\\
Octree & $163\,264$--$377\,856$ & $1\,004\,672$--$2\,691\,008$ & 5,02 / 5,39 & 11,2 / 12,1 & 72,2 / 75,7 & 161 / 171\\\bottomrule
\end{tabularx}
\end{table}

\begin{figure}[H]\centering
\begin{tikzpicture}[x=2.7cm,y=0.40cm,font=\footnotesize]
  \foreach \e/\lab in {-1/{0{,}1},0/{1},1/{10},2/{100}}{
    \draw[muted!50] ({\e+1},-0.3)--({\e+1},-7.7);
    \node[below] at ({\e+1},-7.7) {$\lab$};}
  \node[below] at (3.75,-7.7) {ms};
  \foreach \nome/\v/\lab [count=\i] in {griglia polare/0.40/{0{,}40},lettura nell'immagine/0.53/{0{,}53},colonne di bit/38.4/{38{,}4},griglia densa/49.1/{49{,}1},hash a blocchi/76.4/{76{,}4},tre fasce di quota/97.6/{97{,}6},octree/171/{171}}{
    \pgfmathsetmacro{\x}{log10(\v)+1}
    \fill[accent!65] (0,-\i+0.3) rectangle (\x,-\i-0.3);
    \node[left] at (0,-\i) {\nome};
    \node[left,text=white,inner sep=2pt] at (\x,-\i) {$\lab$};}
  \pgfmathsetmacro{\xt}{log10(66.7)+1}
  \draw[dashed,thick] (\xt,0) -- (\xt,-7.7);
  \node[above] at (\xt,0) {$T=66{,}7$ ms};
\end{tikzpicture}
\caption{95° percentile del tempo per fotogramma con tutti i raggi ($k=1$) e $s=5\unita{cm}$, in scala logaritmica, contro il periodo $T$ (dati della tabella~\ref{tab:cos-tempi}).}\label{fig:cos-tempi}
\end{figure}

Dalle misure seguono quattro constatazioni.
\begin{enumerate}
\item \emph{Il costo segue il prodotto raggi $\times$ passi, non la struttura.} La griglia polare costa 0,03--0,6 ms; la lettura nell'immagine 0,5 ms a fotogramma più 5,9 ns a punto interrogato (636 µs per $108\,000$ punti); la griglia densa 2--4 ms con un raggio su 16 e 45--67 ms con tutti i raggi a 5 cm (44,5 / 49,1 ms sul nucleo veloce; 65,2 / 67,4 ms su quello lento, che in mediana costa 1,49 volte).
\item \emph{La sparsità costa tempo e risparmia poco.} Sulle medie di 20 configurazioni la tabella hash costa 1,39--1,76 volte la griglia densa e l'octree 2,2--3,8 volte (2,3--3,9 con la fusione dei figli uguali, che non riduce il picco dei nodi: almeno il 99\,\% del non fuso). La tabella hash occupa dal 33 all'84\,\% della densa (rapporto 1,2--3,0, con profondità valida fino a 6 m; il limite del lemma~\ref{lem:cos-sparsa} vale a 3 m): il risparmio massimo è $2\,304\,000-763\,904=1\,540\,096$ byte, cioè 1,54 MB. L'octree nella scena ingombra occupa più della densa (117--131\,\%).
\item \emph{Il linguaggio pesa più della struttura.} Lo stesso aggiornamento denso (10 cm, $19\,200$ raggi) scritto in Python senza librerie numeriche costa 102 ms in mediana, contro 2,0 ms in C; la griglia polare 1,6 ms.
\item \emph{Confronto con il periodo.} Con $k=4$ tutte le strutture stanno sotto 12,8 ms al 95° percentile (20,4 ms sul nucleo lento). Con tutti i raggi a 5 cm (figura~\ref{fig:cos-tempi}) la griglia densa occupa il 74\,\% del periodo sul nucleo veloce e lo supera su quello lento (67,4 ms); tabella hash e octree lo superano sempre.
\end{enumerate}
Una ripetizione delle misure sulla stessa macchina, con altri processi in esecuzione, ha dato medie entro il $\pm10\,\%$ in 29 configurazioni su 32. Un secondo programma indipendente (voxel da 8 byte, stanza chiusa, 33,3 milioni di passi) dà 88,9 ms per la griglia densa a 5 cm con tutti i raggi, cioè 2,67 ns per passo, contro 2,6 ns del primo sullo stesso tipo di nucleo.

\subsection{Costo della catena completa}\label{sec:cos-catena}

\textbf{Previsione durante l'occlusione.} Il calcolo di $\Reg_\alpha(t)$ (\S\ref{sec:occ}) per gli oggetti senza osservazione costa, in C, 0,6--1,4 ms al 95° percentile con dieci oggetti nascosti insieme (due programmi, due geometrie dell'ostacolo) e 5,4 ms con 120, cioè l'8,2\,\% di $T$: cresce in modo lineare, circa 0,05 ms a oggetto ($N=1\,500$ cicli, scena sintetica). In Python senza librerie numeriche esistono due misure, entrambe su $N=150$ cicli con dieci oggetti: 5,3 ms di mediana e 6,2 ms al 95° percentile in una prima misura; 6,4 ms di mediana e 15,9 ms al 95° percentile in una ripetizione con altri processi in esecuzione sulla stessa macchina. Il tetto assegnato a questo calcolo è 6,7 ms, un decimo di $T$: la prima misura lo rispetta, la seconda no; il valore in Python non è stabile sotto carico. \prova{S}

\textbf{Catena di riferimento} (mappa polare, gruppi, associazione, filtro di Kalman, regole, uscita; \S\ref{sec:sti}), in Python con una libreria numerica, un solo filo, stesso calcolatore: tabella~\ref{tab:cos-catena}.

\begin{table}[H]\centering\small
\caption{Tempo per ciclo della catena di riferimento (ms). La terza riga non è stata ricalcolata dai dati per ciclo.}\label{tab:cos-catena}
\begin{tabularx}{\textwidth}{@{}Lrrrrc@{}}\toprule
Ingresso & $N$ cicli & Media & 95° perc. & Massimo & Livello\\\midrule
Profondità simulata dalle schede tecniche, $424\times240$, 162 sequenze & $9\,152$ & 6,35 & 18,6 & 105,0 & \prova{S}\\
La stessa, con la soglia minima di punti per ostacolo portata da 3 a 40 & $9\,152$ & 0,95 & 1,28 & 7,7 & \prova{S}\\
Scene senza rumore, $848\times480$ & 48 & 2,5 & 2,9 & 4,7 & \prova{S}\\
Registrazione pubblica affollata, $640\times360$ \cite{depthtrack} & 320 & 81 & 152 & 383 & \prova{P}\\\bottomrule
\end{tabularx}
\end{table}

Sulle sequenze simulate 4 cicli su $9\,152$ superano 66 ms (0,04\,\%; intervallo di Wilson al 95\,\%: $[0{,}02;\,0{,}11]$\,\%). Sulla registrazione pubblica la media supera il periodo: in 21,3 s la catena conferma 254 oggetti, e ne tiene insieme 152 in mediana (le due cifre misurano cose diverse: gli oggetti confermati nell'intera sequenza e quelli presenti nello stesso ciclo).

\textbf{Dipendenza dal numero di oggetti seguiti.} A ogni ciclo la catena confronta ogni oggetto seguito con ogni gruppo (porta $\gamma$ sull'innovazione $\nu$ con covarianza $S$), risolve un'assegnazione ottima \cite{kuhn1955} e aggiorna un filtro per oggetto: lavoro almeno lineare in $n_T$ e polinomiale nell'assegnazione. Sui $9\,152$ cicli simulati il tempo medio cresce con $n_T$: 0,9 ms con 0--2 oggetti, 2,6 ms con 6--10, 5,3 ms con 11--20, 8,1 ms con 21--40, 17,6 ms con 41--80, 34,0 ms oltre 80; ai minimi quadrati
\begin{equation}\label{eq:cos-retta}
t\approx 0{,}33+0{,}30\,n_T\unita{ms}\qquad(\text{termine quadratico } 1{,}7\cdot10^{-4}\,n_T^2).
\end{equation}
Con $n_T=152$ la retta dà circa 46 ms, 50 ms con il termine quadratico, lo stesso ordine degli 81 ms misurati sulla registrazione pubblica; l'estrapolazione oltre i valori di $n_T$ simulati non è stata verificata. Gli stessi fotogrammi, con la soglia a 40 punti (mediana 2 oggetti seguiti invece di 17), costano 0,95 ms. Il costo si governa quindi limitando gli oggetti spuri, non cambiando la rappresentazione dello spazio. \prova{S}

\begin{riquadro}{Limiti di validità delle misure di costo}
(1)~Tutti i tempi sono di un calcolatore portatile, con un solo filo, su profondità sintetica o registrata: la scheda di bordo non è misurata; il solo cambio di tipo di nucleo vale 1,49 volte. (2)~La catena completa nel linguaggio finale non è misurata: le strutture sono misurate in C, la catena in Python, e il divario fra i due linguaggi sull'aggiornamento denso è di circa 50 volte. (3)~La latenza $\tau$ da un estremo all'altro non è misurata (\S\ref{sec:lim}). (4)~Octree e strutture sparse non sono stati misurati con una libreria esistente né su regioni grandi. \prova{N}
\end{riquadro}

\subsection{Alternative considerate ed escluse}\label{sec:cos-alternative}

La tabella~\ref{tab:cos-alternative} elenca le alternative e ciò che le esclude. Dove non è detto altro le prove sono in simulazione, con modello ideale del sensore, camera ferma a $h=1{,}2\unita{m}$ e $\theta=15^\circ$, 15 Hz.

\begin{table}[H]\centering\footnotesize
\caption{Alternative considerate ed escluse. Le frequenze sono date come conteggi; dove la fonte della misura non riporta la numerosità è scritto «$N$ non riportata».}\label{tab:cos-alternative}
\begin{tabularx}{\textwidth}{@{}>{\RaggedRight\arraybackslash}p{3.3cm}>{\RaggedRight\arraybackslash}p{4.4cm}Lc@{}}\toprule
Alternativa & Motivo tecnico & Controesempio o misura & Liv.\\\midrule
Tre sensori a tempo di volo multizona per coprire $360^\circ$ (\S\ref{sec:geo}) & Il campo di un sensore è $60^\circ\times60^\circ$ \cite{st-vl53l7cx}; una zona non distingue il suolo da un ostacolo basso & Tre sensori affiancati coprono $3\cdot60^\circ=180^\circ$: per $360^\circ$ ne servono sei. Una zona è larga 26 cm a 2 m. Portata dichiarata a 5 klx nel modo $8\times8$: 0,5 m \cite{st-vl53l7cx} & \prova{C}\\
Stima della massa (\S\ref{sec:sti}) & Non identificabile da misure di sola posizione con forza ignota & $(M,u)$ e $(\lambda M,\lambda u)$ danno la stessa legge delle misure; informazione di Fisher di rango 1 & \prova{D}\\
Stima dell'accelerazione (\S\ref{sec:sti}) & L'accelerazione fotogramma per fotogramma sta sotto il rumore & $\sigma_a=\sqrt6\,\sigma_p/\Delta^2=16{,}5\unita{m/s^2}$ a 15 Hz con $\sigma_p=3$ cm (valore assunto), contro $0{,}67\unita{m/s^2}$ efficaci nel modello di pedone. Filtro ad accelerazione costante dopo 2 s senza misure: errore 2,29 m contro 1,18 m ($1\,500$ prove per caso, rumore assunto 5 cm) & \prova{C} \prova{S}\\
Catena filtro $\to$ rete neurale $\to$ decisione & Gli enunciati sul «libero» e sull'avviso sono dimostrati, nel modello dichiarato e sotto le sue ipotesi, su regole scritte (\S\ref{sec:map}, \S\ref{sec:occ}) & Scelta di progetto: nessuna misura la esclude. Un componente appreso assegna solo un'etichetta (vedi testo) & \prova{N}\\
Mappa 3D a voxel: densa, sparsa a hash, octree & A sensore fermo non contiene più del fotogramma corrente; le strutture sparse costano tempo & Hash: 1,39--1,76 volte il tempo della densa per al più 1,54 MB; octree: 2,2--3,8 volte. Densa a 5 cm con tutti i raggi: 44,5 / 49,1 ms (tabella~\ref{tab:cos-tempi}) & \prova{S}\\
Varianti che ricordano o deducono un «libero» (\S\ref{sec:map}) & Un «libero» ricordato non è garantito in un mondo che cambia & Sette varianti, ciascuna con il suo controesempio, nella Tabella~\ref{tab:map-varianti}; per esempio la lettura a due stati dà $80\,505$ falsi su $146\,901$ punti dati per liberi (54,8\,\%). Lettura del fotogramma corrente, stesse scene: 0 & \prova{S}\\
Mappa a tre fasce di quota & Parte da 0,15 m; non costa meno della griglia densa & Dichiara il 30,6\,\% del libero vero (0,0\,\% con il 5\,\% di pixel non validi sparsi); 32 celle «libere» sopra una scatola alta 10 cm (griglia polare: 0); con la fascia alta, 0 celle libere su $2\,610$ entro 3 m; 95° percentile 3,4--6,4 ms contro 2,1--4,4 & \prova{S} \prova{C}\\
Previsione a due ipotesi di moto con pesi aggiornati a ogni ciclo, regione a ellisse & I pesi contano la stessa evidenza a ogni ciclo; l'ellisse non è una regione tarata & Sosta dichiarata nel 38--100\,\% dei passaggi normali (cinque simulatori). Regione assente per almeno un ciclo in 374 e $1\,312$ passaggi su $2\,000$, e in $1\,976$ e $1\,987$ su $2\,000$ (due geometrie con ombra larga, due letture della regola). Ellisse al 95\,\%: vedi testo & \prova{S}\\
Scelta automatica della configurazione sotto carico & Ridurre la mappa in esercizio crea «libero» non osservato; il tempo dipende dalla configurazione scelta & «Libero» sul 13,8\,\% e sul 57,0\,\% della zona (due simulatori); 200 cambi in 10 minuti con l'isteresi; 17,5--28 cambi al minuto (quattro simulatori); $N$ non riportata & \prova{S}\\
Probabilità d'intrusione come criterio d'allarme & Scende quando il dato peggiora; non dice se il soggetto \emph{può} entrare & Su rotta di collisione esatta, a 2 s, vale 0,16--0,18; vale 0 con il soggetto nascosto a 9 cm dal volume protetto (\S\ref{sec:occ}) & \prova{D} \prova{S}\\
Sensori laterali come preallarme (\S\ref{sec:geo}) & Il margine temporale guadagnato è di centesimi di secondo & A 1,4 m/s restano 0,07 s; chi passa a 0,5 m di lato entra nel campo a 0,73 m, dentro la zona vicina che la camera non vede. Resta possibile un canale di sola prossimità laterale, che non dichiara mai «libero» & \prova{C}\\
Predittori appresi delle traiettorie & Sull'orizzonte breve il guadagno sul punto è piccolo; la taratura dell'incertezza non è garantita & Nessuna prova in questo lavoro; dati di letteratura \cite{scholler2020,cheraghipouria2026} (vedi testo) & \prova{N}\\\bottomrule
\end{tabularx}
\end{table}

\textbf{Mappe volumetriche.} L'argomento di sicurezza contro una mappa persistente a sensore fermo (Osservazione~\ref{oss:map-fermo}), i controesempi che lo sostengono (Tabella~\ref{tab:map-varianti}) e il conto sulla latenza della griglia bayesiana (Proposizione~\ref{prop:map-logodds}) sono nel \S\ref{sec:map-persistente}; qui restano i conti che riguardano il costo. La decimazione dei raggi non salta voxel solo se $k\le s\,f/d_{\max}$: con $f=337$ pixel e $d_{\max}=6\unita{m}$, $k\le2$ a 5 cm e $k\le5$ a 10 cm \prova{C}. La mappa a tre fasce di quota, con la fascia alta, non dichiara mai libero entro 3 m perché la quota di 2,0 m entra tutta nel campo solo da $z=(2{,}0-h)/\tan(58^\circ/2-\theta)=0{,}8/\tan14^\circ=3{,}21\unita{m}$ \prova{C}. Una mappa persistente tornerebbe necessaria con il sensore in moto, caso non verificato né sperimentalmente né in simulazione (\S\ref{sec:lim}): la forma compatibile con i risultati del \S\ref{sec:map} è una griglia densa locale a tre stati con l'età del dato, che aggiunge avvisi «occupato ricordato» e non produce mai «libero»; eseguita su un altro nucleo, in simulazione porta il 95° percentile della griglia polare da 0,84 a 0,88 ms. Le strutture sparse diventano convenienti quando la regione cresce: $20\times20\times5\unita{m}$ a 5 cm sono $16\cdot10^6$ voxel, 128 MB a 8 byte.

\textbf{Previsione a due ipotesi.} L'alternativa mantiene due ipotesi di moto (transito e sosta), ne moltiplica i pesi a ogni ciclo, dichiara la regione come ellisse al 95\,\% della miscela e smette di mostrarla oltre una soglia fissa di incertezza (0,8 m). In simulazione i pesi moltiplicati a ogni ciclo contano più volte la stessa evidenza e dichiarano una sosta in una quota dei passaggi normali che va dal 38 al 100\,\% secondo il simulatore. La soglia a 0,8 m non scatta mai con un pannello largo 1 m (0 passaggi su $3\,600$) e, in due geometrie con ombra larga, fa mancare la regione nei passaggi della tabella~\ref{tab:cos-alternative}. La prima misura alla ricomparsa cade dentro l'ellisse dichiarata al 95\,\%, per il soggetto che prosegue, in 0,28 / 0,30 / 0,00 dei passaggi con una lettura della regola e in 0,997 / 0,71 / 0,00 con l'altra, in tre geometrie dell'ostacolo (la prima con pannello largo 1 m, la terza con pannello largo 2 m); la dicitura «$\pm\sigma$» corrisponde a una regione che contiene il 39,3\,\%. Alla prova di taratura la regola, applicata alla lettera, copre il 76,4\,\% $[70{,}7;\,81{,}3]$ contro il 90--95\,\% richiesto (nella lettura più favorevole: 100\,\% per chi prosegue, 46,6\,\% per chi torna indietro); anche con i pesi corretti chi esce dal lato lontano dell'ombra è contenuto nel 28\,\% e nel 7\,\% dei casi a 2 e 3 s. Dell'alternativa restano il ritaglio della regione sul suolo visto libero, con la fascia di guardia, e le regole di disegno, fra cui quella per cui la regione non sparisce (\S\ref{sec:occ}). \prova{S}

\textbf{Scelta automatica della configurazione.} Un selettore che in esercizio cambia la configurazione sotto carico e riduce la mappa fa nascere «libero» su suolo non osservato e, se la griglia viene riallocata, perde l'intero insieme di possibile presenza $\Pres(t)$; l'isteresi nel tempo non ferma l'oscillazione quando il tempo di ciclo dipende dalla configurazione scelta: ne fissa solo il ritmo. Nel sistema la configurazione si sceglie e si controlla all'avvio, è dichiarata in ogni uscita, e in esercizio la mappa non cambia (\S\ref{sec:lim}). \prova{S}

\textbf{Rete neurale e classificazione.} Nel sistema decide una catena di rilevamento geometrico, associazione, filtro di Kalman e regole scritte; nessun modello appreso produce un «libero», toglie un avviso o abbassa un livello. Il componente di classificazione su più classi che il sistema richiede (\S\ref{sec:sti}) sta fuori dal percorso che decide: assegna un'etichetta all'oggetto seguito, in modo asincrono, e l'etichetta può solo allargare l'avviso (\S\ref{sec:occ}). In questo lavoro il classificatore non è realizzato né valutato: è simulato con etichette assegnate a mano (\S\ref{sec:sti}). Un conto mostra che per le classi veloci l'etichetta arriverebbe comunque tardi: i 3 m di portata si attraversano in 0,75 s a 4 m/s, e la sola conferma di un oggetto seguito chiede 4 cicli, cioè 0,27 s; per queste classi protegge la configurazione iniziale, non l'etichetta \prova{C}. Nessun rilevatore di oggetti è stato eseguito in questo lavoro. Dalla documentazione pubblica, i rilevatori pronti addestrati sull'insieme pubblico di immagini COCO coprono, secondo un conto fatto sull'elenco delle sue classi, 9 dei circa 25 tipi di una tassonomia più fine di quella della Tabella~\ref{tab:sti-profili}, che distingue per esempio la persona su carrozzina (mancano carrello elevatore, tram, multirotore, carrozzina), e chiedono l'immagine a colori; la letteratura riporta che le reti moderne danno probabilità mal tarate e propone una correzione a posteriori. Queste affermazioni di letteratura non sono state riscontrate su fonti presenti in bibliografia e sono date senza riferimento. Se un modello di questo tipo basti per queste classi, a 3 m e sulla sola profondità, non è verificato. \prova{N}

\textbf{Predittori appresi delle traiettorie.} Sull'orizzonte di 4,8 s, sulla media dei cinque scenari degli insiemi pubblici ETH e UCY \cite{pellegrini2009,lerner2007}, il modello a velocità costante ha errore medio e finale di 0,39 m e 0,83 m, contro 0,44 m e 0,91 m del miglior predittore appreso deterministico e 0,32 m e 0,60 m di un modello generativo valutato sul migliore di 20 campioni \cite{scholler2020}; i predittori addestrati sulla verosimiglianza tendono a distribuzioni troppo o troppo poco sicure \cite{cheraghipouria2026}. Il confronto riguarda il punto previsto, non una regione, non comprende occlusioni, e quei dati non vengono da un sensore di profondità frontale. Per questo il sistema usa il modello a velocità quasi costante e dichiara la regione con un livello misurato a parte (\S\ref{sec:val}); nessun predittore appreso è stato provato in questo lavoro.

\subsection{Lavori correlati}\label{sec:cos-correlati}

\textbf{Sensori di profondità.} Per le camere a stereo attivo la documentazione del costruttore dichiara che l'errore di profondità cresce con il quadrato della distanza \cite{realsense-taratura-profondita}; una caratterizzazione indipendente su un bersaglio planare bianco fra 0,4 e 2,1 m dà per il rumore assiale un polinomio di secondo grado nella distanza \cite{ahn2019}. I sensori a tempo di volo multizona offrono al più $8\times8$ zone a 15 Hz, oppure $4\times4$ zone a 60 Hz \cite{st-um2884}, con una portata che la luce ambiente riduce \cite{st-vl53l7cx}; l'uso all'aperto delle camere a tempo di volo resta limitato dall'interferenza solare \cite{qiu2023}; un LiDAR 2D rotante economico misura un solo strato orizzontale, a circa 10 Hz e con risoluzione angolare di $0{,}72^\circ$ \cite{slamtec-rplidar-c1}. Sono dichiarazioni dei costruttori o misure di terzi, non misure fatte sul sistema descritto; il loro uso nei calcoli è nel \S\ref{sec:geo}.

\textbf{Mappe di occupazione e stato «sconosciuto».} Le griglie di occupazione sono la rappresentazione classica per la percezione dei robot mobili \cite{elfes1989}. OctoMap rappresenta in modo esplicito, oltre allo spazio occupato, anche quello libero e quello sconosciuto \cite{hornung2013}; nella sua realizzazione lo sconosciuto è codificato in modo implicito, come assenza di informazione \cite{octomap}. UFOMap dichiara una rappresentazione esplicita di tutti e tre gli stati \cite{duberg2020}. Le griglie evidenziali trattano l'ignoranza come una sorgente d'incertezza distinta dall'errore di misura \cite{moras2011}, e le griglie dinamiche di occupazione sono state realizzate in tempo reale con un filtro a particelle in parallelo \cite{nuss2018}. Le mappe volumetriche a distanza con segno funzionano in tempo reale su un solo nucleo di CPU \cite{oleynikova2017} o su GPU \cite{millane2024,niessner2013}, e per i veicoli in volo è stata proposta una griglia locale, in un buffer circolare, che si muove con il veicolo \cite{usenko2017}. Rispetto a questi lavori il sistema adotta una mappa polare del suolo, centrata su una camera ferma, con quattro stati: separa lo sconosciuto in «in ombra» $\Omb(t)$ (dietro un ostacolo osservato) e «non osservato», e non aggiorna una probabilità per cella. Non è stata confrontata sperimentalmente con le librerie citate: i confronti del \S\ref{sec:cos-misure} usano realizzazioni proprie delle stesse strutture.

\textbf{Seguimento di più oggetti e occlusioni.} Lo schema «rilevamento, filtro di Kalman \cite{kalman1960}, assegnazione ottima \cite{kuhn1955}» è quello di SORT, che dato l'insieme dei rilevamenti aggiorna gli oggetti seguiti a 260 Hz \cite{bewley2016}, e della sua estensione in tre dimensioni \cite{weng2020}; la teoria della stima per il seguimento è materia di manuale \cite{barshalom2001}. I perfezionamenti noti agiscono sull'associazione (uso di quasi tutti i rilevamenti, anche a basso punteggio \cite{zhang2022}; compensazione del moto della camera \cite{aharon2022}) e sulla gestione del ciclo di vita, isolata come modulo a sé \cite{pang2021}. Durante un'occlusione il filtro, privo di misure, accumula errore \cite{cao2023}: è il punto in cui si colloca questo lavoro. La letteratura segue tre strade. La prima tratta l'occlusione come un problema di previsione a breve termine \cite{khurana2021} o la impara da dati che portano le annotazioni degli oggetti invisibili \cite{tokmakov2021}; i lavori di riferimento riconoscono però che i metodi faticano sotto occlusione pesante \cite{hsieh2023} e che resta un divario notevole rispetto a una vera permanenza dell'oggetto \cite{vanhoorick2023}. La seconda usa l'informazione «negativa», cioè la mancata osservazione là dove il sensore vede \cite{koch2007,wyffels2015,krejci2026}. La terza, nata nella guida automatica, sovrastima con insiemi lo stato di ciò che è nascosto \cite{orzechowski2018} e predice in forma insiemistica i comportamenti ammissibili dei partecipanti rilevati e di quelli potenzialmente nascosti \cite{koschi2021}; il problema affine dell'inseguimento con visibilità considera un evasore imprevedibile, di posizione iniziale ignota e senza limite di velocità \cite{guibas1997}. Il sistema combina la prima e la terza strada: per l'oggetto già seguito mantiene una regione di previsione $\Reg_\alpha(t)$ dichiarata a un livello $\alpha$; per l'ombra in cui è entrato un oggetto seguito mantiene un'area d'avviso $W(t)$, che realizza l'insieme di possibile presenza $\Pres(t)$ del \S\ref{sec:occ} con una velocità massima assunta $\vmax$; un'ombra in cui non è entrato alcun oggetto seguito resta «in ombra», senza avviso. La percezione umana offre un criterio per la resa: l'identità di un oggetto nascosto è mantenuta se scompare e ricompare lungo il bordo fisso dell'ostacolo \cite{scholl1999}.

\textbf{Incertezza tarata.} Le previsioni probabilistiche vanno valutate per affidabilità e nitidezza \cite{kress2021} con punteggi propri \cite{gneiting2007}; un'analisi recente conferma che la velocità costante approssima bene la dinamica dei pedoni nella maggior parte dei casi \cite{uhlemann2024}. La previsione conforme dà regioni con copertura marginale garantita senza ipotesi sulla distribuzione, purché i dati di taratura e di prova siano scambiabili \cite{angelopoulos2021}; è stata applicata alla pianificazione in ambienti dinamici, con regioni su 20 passi e 498 traiettorie di prova su 500 dentro la regione al 95\,\% \cite{lindemann2023}, resa meno prudente \cite{cleaveland2024} e adattata in linea quando la distribuzione cambia \cite{gibbs2021,dixit2022}. Due limiti noti valgono anche qui: la garanzia è marginale, e una copertura condizionata esatta è impossibile senza ipotesi sulla distribuzione \cite{barber2019}; una probabilità di collisione calcolata con dati peggiori può apparire più bassa \cite{balch2019}, il che motiva l'uso di regioni e insiemi anziché di una sola probabilità.

\textbf{Sicurezza dei sensori di presenza.} Nella sorveglianza di velocità e separazione dei robot collaborativi la distanza di protezione ha forma lineare, somma dei percorsi dell'operatore e del robot nei tempi di reazione e di arresto e di termini d'incertezza di posizione, con una velocità dell'operatore di 1,6 m/s \cite{marvel2017}. I sistemi 3D industriali di questo tipo classificano i volumi in vuoti, occupati e sconosciuti e trattano come occupato ciò che nessun sensore vede \cite{therobotreport-veo,us-brevetto-11846916}; il segnalamento ferroviario applica lo stesso principio a guasto sicuro, imponendo l'aspetto più restrittivo quando manca l'energia \cite{us-49cfr-236-205}. Il sistema ne riprende due principi: non dichiarare mai libero ciò che non è osservato, e segnalare in modo esplicito la perdita di affidabilità. Non è un dispositivo di protezione certificato, e nessuna conformità a norme di sicurezza delle macchine è rivendicata; le norme sulla distanza di separazione non sono state consultate nel testo originale.

\textbf{Ausili per persone con disabilità visiva.} Le rassegne mostrano che efficacia e sicurezza sono la prima preoccupazione degli utenti \cite{xu2023} e che solo il 22,9\,\% di 70 lavori valuta il dispositivo con persone con disabilità visiva \cite{romlay2021}. Un sistema di avviso con previsione per estrapolazione lineare e soglie a 5 s e 2,5 s è stato provato con sei utenti non vedenti \cite{kayukawa2019}; esistono dispositivi indossabili con sensori di profondità e riscontro aptico \cite{dossantos2025,kozlov2026}. Sul canale d'uscita gli studi sono discordi: per chi usa il bastone il riscontro sonoro ha dato un carico cognitivo 22 volte maggiore di quello aptico, effetto assente nei partecipanti bendati \cite{martinez2014}; in un altro studio, con otto partecipanti, gli stimoli aptici sono risultati meno intuitivi di quelli sonori \cite{lupu2020}; la guida con suono spazializzato regge un compito secondario meglio del linguaggio spaziale \cite{giudice2008}. Il sistema non risolve questa scelta: espone gli stati e l'incertezza in forma indipendente dal canale (\S\ref{sec:lim}).
